\fontsize{13pt}{22.5pt}\selectfont
\begin{center}
    {\large\bfseries TÓM TẮT ĐỒ ÁN}
\end{center}

\vspace{0.5cm}

\indent Trong một dự án học tăng cường cho game, phần việc nặng nhất hiện không còn nằm ở
thuật toán. Các thư viện đã đóng gói sẵn những gì cần thiết. Cái còn lại, và cũng là cái
khó, là môi trường: phải tự dựng, tự đặc tả, tự hiệu chỉnh. Đồ án này lấy chính khâu đó
làm đối tượng nghiên cứu, cụ thể là bài toán xây dựng môi trường game ba chiều đủ chuẩn
cho việc huấn luyện.

\indent Ba môi trường đã được dựng trên nền tảng Unity ML-Agents, mỗi cái ứng với một
dạng bài toán khác nhau. Football Table là bài toán đối kháng một-một với không gian hành
động liên tục và cơ chế tự chơi. Capture The Flag đòi hỏi hai tác nhân hợp tác, dùng phần
thưởng nhóm cùng chương trình học chia theo bậc. Cross The Road thì ngược lại, chỉ một tác
nhân điều hướng nhưng phần thưởng rất thưa. Sản phẩm cuối được đóng thành một gói
Unity Package Manager, kèm 15 cấu hình huấn luyện, bộ công cụ Editor và trình hướng dẫn
cài đặt.

\indent Đồ án đã cho chạy năm thuật toán PPO, SAC, MA-POCA, MAPPO và DQN trên cả ba
môi trường với cùng ngân sách bước trong từng môi trường; các siêu tham số đặc thù được
điều chỉnh theo từng thuật toán. Mười lăm lần chạy cho thấy cả ba môi trường đều có thể
dạy được hành vi mục tiêu và không bộc lộ lối tắt khi quan sát trực tiếp. Ở các lần chạy
đã ghi nhận, SAC đạt $0{,}846$ trên Cross The Road nhưng chỉ đạt $0{,}183$ trên Capture
The Flag. Tuy nhiên, mỗi tổ hợp mới được chạy với một seed, nên các chênh lệch này là kết
quả thăm dò chứ chưa phải bằng chứng thống kê để xếp hạng thuật toán. Riêng DQN trên
Football Table được trình bày như tham chiếu độc lập vì không dùng self-play. Kèm theo đó,
đồ án xây dựng một ứng dụng trình diễn nạp trực tiếp 15 mô hình ONNX đã huấn luyện để quan
sát hành vi của các tác nhân.

